%!TEX root = ../main.tex
% =============================================================================
%  第 1 章  问题重述
%  评审在问：你知道题目在问什么吗？
%  本章只做两件事：① 把题面语言改写成自己的语言；② 把“要交什么”列清楚。
%    1.1 问题背景（3–5 句）  1.2 需要解决的问题（编号 4 条，与第 6–9 章一一对应）
%    1.3 交付物（结果文件 ↔ 题面表 1--表 4）
%  约定：记号、决策矛盾、依赖关系与统一建模框架全部放在第 2 章（问题分析）；
%        本章不出现任何模型公式与本题结果数字。
% =============================================================================

\section{问题重述}\label{chap:restate}

\subsection{问题背景}

微网是为解决光伏发电等分布式能源本地消纳、减少并网问题而发展起来的小型发用电系统。在分时电价下，
光伏盈余多出现在电价低谷的午间，而小区负载高峰落在电价尖峰的傍晚，两者在时间上错配；储能可在电价
较低或分布式能源有剩余的时间段充电、在电价较高的时间段放电，因而是微网降低购电费用的主要手段。
本题要求在给定小区负载、光伏发电与电价数据的条件下，逐时段给出微网的购电量与储能充放电量，
使微网提供的电力既满足小区负载、又尽可能节省购电费用，并逐层加入运行不确定性、多时刻预报与
波动电价。四问的差别不在模型复杂度，而在\textbf{信息结构}：决策时刻由 1 个增至 4 个、计划依据由
确定值变为预报值，并引入两类方向相反的偏差惩罚（详见 \S\ref{sec:analysis-depend}）。

\subsection{需要解决的问题}

\begin{enumerate}
  \item \textbf{问题 1（典型日确定性计划）}：设每天的电价与小区负载相同，且已给出光伏发电功率一天的
    预测数据；要求微网提供的电能不低于小区负载，且储能设备 \(0{:}00\) 与 \(24{:}00\) 的储电量相同。
    建立每天 \(0{:}00\) 制定当天计划购电策略的数学模型，并以表 1 格式给出指定时段的购电量、
    全天购电量与全天购电费，以表 2 格式给出指定时段的充放电量与 \(0{:}00\)、\(24{:}00\) 的储电量，
    完整计划购电策略写入 \texttt{result1.xlsx}。
  \item \textbf{问题 2（负载与光伏逐日变化下的日前计划）}：电价每天相同，而小区负载与光伏发电功率
    随时间变化；若微网提供的电能低于负载，须向外部电网\textbf{紧急购电}，电价为交易时刻电价的
    \(5\) 倍，其余时段的购电费用按计划购电量计算。在每天 \(0{:}00\) 制定当天计划购电策略；
    按表 1、表 2 格式给出表 3 中指定日期的结果，按表 3 格式给出紧急购电结果，
    完整策略写入 \texttt{result2.xlsx}。
  \item \textbf{问题 3（多时刻预报下的调整策略）}：每天 \(0{:}00\)、\(6{:}00\)、\(12{:}00\)、\(18{:}00\)
    可获得未来 \(24\)~h 整点的光伏发电功率预报；可据 \(0{:}00\) 预报制定当天计划，并据其他时刻的预报
    调整购电策略。计划购电量高于调整购电量的部分，违约电价为交易时刻电价的 \(50\%\)；
    调整购电量高于计划购电量的部分，超出部分电价为交易时刻电价的 \(150\%\)。总购电费用含计划购电费、
    紧急购电费与调整购电量的相关费用。建立当天购电策略的数学模型，按表 1、表 2、表 3 格式给出
    表 3 中指定日期的结果，完整策略写入 \texttt{result3.xlsx}，并\textbf{分析是否需要引入其他时刻的
    预报}来确定调整购电策略。
  \item \textbf{问题 4（波动电价下的重算）}：实际运行中外网电价亦实时波动。针对波动电价建立微网当天
    购电策略的数学模型，并在该电价数据下重新计算问题 2 与问题 3，完整购电策略分别写入
    \texttt{result4-2.xlsx}（对应问题 2）与 \texttt{result4-3.xlsx}（对应问题 3）。
\end{enumerate}

\subsection{交付物}

五个结果文件与题面指定表格的对应关系见表~\ref{tab:deliverable}。结果文件的工作表名、列序与日期范围
（2025-02-01 至 12-31 共 334 天）与题目附件 5 的模板逐项一致；支撑材料（全部源程序与中间结果）
清单见附录 A。

\begin{table}[htbp]
  \centering
  \caption{交付物与题面指定表格的对应关系（表号指题面的表 1--表 4；\textbf{表 3} 为题目指定的四个日期：
            2025-03-20、06-21、09-23、12-21）}
  \label{tab:deliverable}
  \footnotesize
  \setlength{\tabcolsep}{4pt}
  \begin{tabularx}{\linewidth}{C{2.05cm} C{1.55cm} L{2.55cm} Y}
    \toprule
    结果文件 & 对应问题 & 结果文件交付要求 & 论文中对应表格 \\
    \midrule
    \texttt{result1.xlsx} & 问题 1 & 计划购电量、充放电量（典型日全天 144 槽）
      & 表~\ref{tab:q1-buy}（表 1 格式）、表~\ref{tab:q1-soc}（表 2 格式） \\
    \texttt{result2.xlsx} & 问题 2 & 计划购电量、充放电量、紧急购电量（334 天）
      & 表~\ref{tab:q2-t1}、表~\ref{tab:q2-t2}（四个指定日期）、表~\ref{tab:emg-q2}（表 3 格式） \\
    \texttt{result3.xlsx} & 问题 3 & 计划购电量、调整购电量、充放电量、紧急购电量（334 天）
      & 表~\ref{tab:q3-t1}、表~\ref{tab:q3-t2}（四个指定日期）、表~\ref{tab:emg-q3}（表 3 格式） \\
    \texttt{result4-2.xlsx} & 问题 4（对应问题 2） & 波动电价下的计划购电量、充放电量、紧急购电量
      & 表~\ref{tab:q4-summary}、表~\ref{tab:emg-q42}（表 3 格式） \\
    \texttt{result4-3.xlsx} & 问题 4（对应问题 3） & 波动电价下的计划、调整、充放电与紧急购电量
      & 表~\ref{tab:q4-summary}、表~\ref{tab:emg-q43}（表 3 格式） \\
    \bottomrule
  \end{tabularx}
\end{table}
